\documentclass{article}

\usepackage{amsmath}

\title{Ch. 26 HW \#1}
\author{Principles of Macroeconomics}
\date{\today}

\begin{document}

\maketitle

\section*{Quick Quiz}

\begin{enumerate}

\setcounter{enumi}{5}

\item (c) -- In the long run, the gains from additional capital will eventually fade out until the growth rate has returned to its original level.

\item (c) -- You don't have to know what GNP is, but hopefully you do from reading the textbook. GDP will rise because production on domestic soil increases. GNP -- the income accruing to U.S. citizens regardless of where they live -- will also increase because the factory will create new jobs. GNP will increase less than GDP because much of the income generated from the production will flow back to Japan rather than U.S. citizens. 

\item (a) -- See the ``Stretching Natural Resources" section of 26.3.i.

\end{enumerate}

\section*{Questions for Review}

\begin{enumerate}

\setcounter{enumi}{3}

\item Higher saving leads to higher investment. Higher investment leads to greater levels of physical capital per worker. When workers have more tools, they are more productive, so greater levels of physical capital per worker increases productivity. How much a society produces dictates its income.

A policymaker might not want to raise the saving rate because it would require lower consumption in the present. 

\item Temporarily -- in the long run, the gains from additional capital fade out due to the diminishing marginal product of capital. 

\item Increased trade allows countries to specialize in producing goods and services for which they hold a comparative advantage. That allows them to develop technological knowledge which promotes productivity.

\item It's ambiguous. On the one hand, it will spread the capital stock and the natural resources across a greater number of people which should hurt productivity. On the other hand, having more people makes technological progress more likely which should help productivity.

\item It funds research at colleges and universities. It incentivizes private innovation by granting companies patents that allow them to make money off of their inventions.

\end{enumerate}

\section*{Problems and Applications}

\begin{enumerate}

\item Importing substantial amounts of goods and services is not mutually exclusive with producing large amounts of goods and services. The U.S. imports a substantial amount, but it also produces a substantial amount. Its high level of production determines its standard of living.

\end{enumerate}

\end{document}